% ===================================================================
%  جدول نمبر 3 — لڑکپن اور جوانی۔
%
%  Traduction, faite en 2026, en ourdou d'aujourd'hui, sur l'ido de
%  texto/io. Regles de la colonne : voir 10-jadval-01.tex. Deux
%  scenes, deux numerotations : les cles portent c1 et c2. Ecrit avec
%  outils/ordo.py sous les yeux.
%
%  DIX-SEPT RETOURNEMENTS SUR VINGT-CINQ PAIRES — le meme compte que
%  les colonnes coreenne et tamoule au meme tableau. C'est un DEFILE,
%  et un defile se decrit par attributs : la nourrice AVEC sa coiffe,
%  le mendiant APPUYE sur sa bequille, la vieille APPUYEE sur son
%  baton, le meunier AVEC son chapeau, les abeilles DE la ruche, la
%  fleche DU clocher, le clocher QUI couronne l'eglise. Trois langues
%  a tete finale, trois familles, le meme partage : parigi.py ne lit
%  que l'ido, et les trois motifs d'exemption sont des proprietes de
%  l'ido.
%
%  LE FAC-SIMILE DONNE DEUX FOIS LE RENVOI (18) ET DEUX FOIS LE (19)
%  au bloc c1-09-1 : les arbres fruitiers puis les pechers pour le
%  premier, le presbytere puis les abricotiers pour le second. C'est
%  une faute de l'edition ido, et la source ne bouge pas.
%
%  LE VOCABULAIRE CHRETIEN EXISTE EN OURDOU DEPUIS QUATRE SIECLES, ET
%  IL EST PERSAN. L'Eglise du Pendjab et du Sind a sa langue
%  liturgique en ourdou depuis les missions du dix-septieme siecle :
%  کلیسا l'eglise, پادری le cure, بپتسمہ le bapteme, گاڈ فادر et گاڈ
%  مدر — non, ceux-la s'empruntent — ; la colonne ecrit منہ بولا باپ
%  et منہ بولی ماں, qui sont les mots de l'usage. On ne forge pas de
%  calque persan pour une charge que la langue nomme autrement.
%
%  LE MOT DEVENU UNE INSULTE, CINQUIEME FOIS DU RELEVE. La baraque du
%  bloc c2-14-1 montre « la giganto e la nano ». L'ourdou a بونا, qui
%  traduit exactement le second et sert aujourd'hui d'injure ; la
%  colonne ecrit بہت چھوٹا آدمی, et — pour que la paire reste une
%  paire — بہت بڑا آدمی pour le geant, dont les mots disponibles
%  (دیو, غول) designent des demons. Meme arbitrage que pour le 난쟁이
%  coreen, le குள்ளன் tamoul et le కసాయి telougou.
%
%  LE « cigano » (41) N'EST PAS DANS CE TABLEAU-CI mais au 9 ; en
%  revanche la question du peuple nomme comme un metier revient au
%  bloc c1-23-1 avec le meunier, et elle ne se pose pas : چکی والا est
%  un metier et rien d'autre.
%
%  CE QUI S'EMPRUNTE ICI S'EMPRUNTE PARCE QUE LA CHOSE NE POUSSE PAS
%  AU PENDJAB : چنار le platane — non, celui-la pousse au Cachemire
%  et la langue l'a —, پاپلر le peuplier, آڑو le pecher — celui-la
%  aussi pousse —, خوبانی l'abricot, ٹیولپ la tulipe, کارنیشن
%  l'oeillet. La liste est plus courte qu'au tamoul et qu'au
%  telougou, et la raison est la latitude : le Pendjab a des hivers.
%
%  ET « جولائی » EST EN GRAS AU BLOC c2-01-1. La ligne « ur » a ete
%  posee dans gravuri/korekti.json AVANT que ce tableau soit ecrit,
%  comme pour le tamoul.
% ===================================================================

%%K t03-apar-1 apar
\VUsaut{3.0mm}
\VUtitre{15.0pt}{60}{جدول نمبر 3}
\VUsaut{1.6mm}
\VUtitre{11.4pt}{0}{لڑکپن اور جوانی۔}
\VUsaut{3.4mm}

%%K t03-apar-2 apar
(تیسرا اور چوتھا جدول اِس طرح جوڑے گئے ہیں کہ چار \VUgras{گھریلو
مناظر} میں \VUgras{موسم}، \VUgras{عمر}، \VUgras{خاندان}،
\VUgras{لباس} اور \VUgras{حالت} کے بنیادی تصور دکھائیں۔)

\VUsaut{4.0mm}
%%K t03-tit-1 sub
\VUcentre{12.0pt}{0}{\textit{پہلا منظر۔}}
\VUsaut{3.0mm}

%%K t03-tit-2 sub
\VUpk{10.2pt}{40}{گاؤں میں بپتسمے کی تقریب۔}
\VUsaut{2.4mm}

%%K t03-tit-3 sub
\VUpk{10.2pt}{40}{بہار۔}
\VUsaut{3.0mm}

%%K t03-c1-01-1 p
1. --- ہم \VUgras{بہار} میں ہیں ؛ \VUgras{مارچ}، \VUgras{اپریل} یا
\VUgras{مئی} کے مہینوں میں، \VUgras{ہفتے} کے دنوں میں —
\VUgras{پیر}، \VUgras{منگل} یا \VUgras{بدھ}۔ صبح کے \VUgras{دس بجے}
ہیں۔

%%K t03-c1-02-1 p
2. --- \VUgras{موسم} خوشگوار ہے، \VUgras{ہوا} صاف ہے۔ سورج چمک رہا
ہے۔ کچھ \VUgras{بادل کے ٹکڑے} \textsuperscript{(2)} نیلے
\VUgras{آسمان} \textsuperscript{(1)} میں اٹھ رہے ہیں۔

%%K t03-c1-03-1 p
3. --- \VUgras{درختوں} پر ابھی مشکل سے \VUgras{پتے} ہیں۔ دبلے
\VUgras{پاپلر} \textsuperscript{(3)} اور اونچے \VUgras{چنار}
\textsuperscript{(17)} پر اب تک صرف کلیاں ہیں۔

%%K t03-c1-04-1 p
4. --- \VUgras{ابابیلیں} \textsuperscript{(4)} لوٹ آئی ہیں اور خوشی سے
چہچہاتی ہوئی \VUgras{نوکدار مینار} \textsuperscript{(7)} کے گرد اڑ رہی
ہیں ؛ وہ مینار \VUgras{گھنٹہ گھر} \textsuperscript{(5)} کی چوٹی پر ہے،
اور وہ گھنٹہ گھر گاؤں کے \VUgras{کلیسا} \textsuperscript{(6)} کو تاج
پہناتا ہے۔

%%K t03-tit-4 sub
\VUsaut{3.4mm}
\VUpk{10.2pt}{40}{کلیسا۔}
\VUsaut{3.0mm}

%%K t03-c1-05-1 p
5. --- وہ کلیسا بہت سادہ ہے ؛ اُس کا بڑا \VUgras{دروازہ}
\textsuperscript{(13)} اور موٹی \VUgras{گھڑی} \textsuperscript{(8)}
ہے۔ چھوٹی \VUgras{سوئی} \textsuperscript{(9)} X کے ہندسے پر ہے ؛ وہ
\VUgras{گھنٹے} بتاتی ہے۔ \VUgras{بڑی سوئی} \textsuperscript{(10)} XII
پر ہے (دوپہر) ؛ وہ \VUgras{منٹ} بتاتی ہے۔

%%K t03-c1-06-1 p
6. --- گھنٹہ گھر میں \VUgras{گھنٹے} \textsuperscript{(11)} خوشی سے بج
رہے ہیں۔ چھت کی چوٹی پر پتھر کی ایک چھوٹی \VUgras{صلیب}
\textsuperscript{(12)} کھڑی ہے۔

%%K t03-c1-07-1 p
7. --- کلیسا کا صرف بڑا \VUgras{دروازہ} \textsuperscript{(13)} ہی
نہیں، ایک چھوٹا پہلو دروازہ بھی ہے۔ اُسی چھوٹے دروازے سے، اپنا
\VUgras{چوغہ} پہنے ہوئے \VUgras{پادری صاحب} \textsuperscript{(15)}
\VUgras{لباس خانے} \textsuperscript{(14)} سے نکل کر \VUgras{پادری کے
گھر} \textsuperscript{(19)} کی طرف جاتے ہیں۔

%%K t03-c1-08-1 p
8. --- اُن کے قریب، یہ رہی \VUgras{دودھ والی} \textsuperscript{(24)} ؛
اُس کے ساتھ اُس کا \VUgras{گدھا} \textsuperscript{(23)} ہے۔ وہ شہر سے
لوٹ رہی ہے ؛ وہاں بچوں کے لیے اچھا دودھ لے گئی تھی۔

%%K t03-tit-5 sub
\VUsaut{4.0mm}
\VUpk{10.2pt}{40}{پھلوں کا باغ۔}
\VUsaut{3.4mm}

%%K t03-c1-09-1 p
9. --- الیبورون صاحب (گدھا) اُس صبح کی سیر سے پوری طرح مطمئن ہے۔ وہ
مہکتی ہوا مزے سے کھینچتا ہے ؛ وہ مہک \VUgras{پھولوں}
\textsuperscript{(20)} کی ہے، اور وہ پھول ساتھ والے پھلوں کے باغ کے
\VUgras{پھل دار درختوں} \textsuperscript{(18)} کو ڈھانپے ہوئے ہیں۔ وہ
\VUgras{آڑو کے درخت} \textsuperscript{(18)} اور \VUgras{خوبانی کے
درخت} \textsuperscript{(19)} بالکل گلابی اور سفید ہیں ؛ اچھی فصل کی
امید دلاتے ہیں۔

%%K t03-c1-10-1 p
10. --- اُسی وقت چھوٹی \VUgras{شہد کی مکھیاں} \textsuperscript{(22)}
وہاں لوٹ مارنے جا رہی ہیں ؛ وہ \VUgras{چھتے} \textsuperscript{(21)}
کی ہیں، اور بھنبھناتی ہوئی اپنا میٹھا \VUgras{شہد} جمع کرتی ہیں۔

%%K t03-c1-11-1 p
11. --- مگر کیا ہو رہا ہے؟ یہ رہے گاؤں کے لڑکے، دوڑ کر آتے ہیں اور
ڈری ہوئی \VUgras{ہنسوں} \textsuperscript{(25)} کو بھگاتے ہیں۔ یہ تو
نوٹری کے بیٹے کے \VUgras{بپتسمے} کے بعد کلیسا سے مجمع کا نکلنا ہے۔

%%K t03-c1-12-1 p
12. --- ننھا یاکوبس اپنے \VUgras{جھولے} \textsuperscript{(26)} میں
گھنٹوں کی خوشی بھری آواز سے جاگ اٹھتا ہے ؛ اور اُس کی بہن مگدلینا،
جو خود بھی بپتسمے کا جلوس دیکھنا چاہتی ہے، اپنی ماں سے کہتی ہے کہ
گھر کی دہلیز پر آ کر اُس کے بال بنا دے اور اُسے کپڑے پہنا دے۔

%%K t03-c1-13-1 p
13. --- ماں مان جاتی ہے۔ وہ اپنا \VUgras{دوپٹہ} \textsuperscript{(28)}،
اپنی \VUgras{انگیا} \textsuperscript{(29)} اور اپنا \VUgras{تہبند}
\textsuperscript{(30)} پہن کر دروازے کے سامنے ایک چھوٹی کرسی پر آ
بیٹھتی ہے۔ مگدلینا کے پاس صرف اُس کا \VUgras{پیٹی کوٹ}
\textsuperscript{(31)} اور اُس کی \VUgras{چولی} \textsuperscript{(32)}
ہیں۔

%%K t03-c1-14-1 p
14. --- وہ کتنی خوش ہے! اپنے پسندیدہ پھول بھول جاتی ہے — اپنے
\VUgras{کارنیشن} \textsuperscript{(34)}، جن پر \VUgras{تتلیاں}
\textsuperscript{(35)} آ بیٹھتی ہیں ؛ اپنے \VUgras{ٹیولپ}
\textsuperscript{(36)} ؛ اپنے \VUgras{مئی کے پھول}
\textsuperscript{(37)} — \VUgras{گملوں} \textsuperscript{(38)} میں،
\VUgras{دیوار} \textsuperscript{(33)} پر — ؛ اپنے \VUgras{گلاب}
\textsuperscript{(39)}، جو اتنی خوبصورت باڑ بناتے ہیں۔ وہ جلوس کی
عورتوں کے قیمتی لباس دیکھ رہی ہے۔

%%K t03-c1-15-1 p
15. --- گاؤں کے لڑکے لباس پر زیادہ دھیان نہیں دیتے۔
\VUgras{مٹھائیاں} \textsuperscript{(55)} یا \VUgras{ریزگاری}
\textsuperscript{(52)} پانے کے لیے ایک دوسرے کو دھکیلتے ہوئے دوڑتے
ہیں۔ اُن میں سے ایک رو رہا ہے \textsuperscript{(40)}۔

%%K t03-c1-16-1 p
16. --- دوسرے نے \VUgras{کتے} \textsuperscript{(42)} کے پاؤں پر پاؤں
رکھ دیا ؛ وہ چیختا ہوا بھاگتا ہے اور \VUgras{مرغی}
\textsuperscript{(43)} اور اُس کے \VUgras{چوزوں} \textsuperscript{(44)}
کو بھگا دیتا ہے۔

%%K t03-tit-6 sub
\VUsaut{5.0mm}
\VUpk{10.2pt}{40}{بپتسمے کا جلوس۔}
\VUsaut{4.0mm}

%%K t03-c1-17-1 p
17. --- جلوس کے آگے \VUgras{دایہ} \textsuperscript{(45)} چل رہی ہے۔
اُس کے سر پر لمبے ربنوں والی خوبصورت \VUgras{اوڑھنی}
\textsuperscript{(46)} ہے۔ وہ \VUgras{نومولود} \textsuperscript{(47)}
کو اٹھائے ہوئے ہے ؛ وہ بچہ سر تا پا \VUgras{جالی دار کام}
\textsuperscript{(48)} میں لپٹا ہوا ہے۔

%%K t03-c1-18-1 p
18. --- دایہ کے پیچھے \VUgras{منہ بولا باپ} \textsuperscript{(49)} اور
\VUgras{منہ بولی ماں} \textsuperscript{(53)} آتے ہیں۔ وہ مٹھائیاں اور
ریزگاری بانٹتے ہیں۔ منہ بولے باپ کے پاس \VUgras{دستانے}
\textsuperscript{(51)} اور \VUgras{اونچی ٹوپی} \textsuperscript{(50)}
ہیں۔ منہ بولی ماں ہاتھ میں خوبصورت \VUgras{گلدستہ}
\textsuperscript{(54)} پکڑے ہوئے ہے۔

%%K t03-c1-19-1 p
19. --- اُن کے پیچھے ہمیں بچے کے \VUgras{چچا} \textsuperscript{(56)}
اور \VUgras{چچی} \textsuperscript{(58)} نظر آتے ہیں۔ چچی کے سر پر
نفیس \VUgras{ٹوپی} \textsuperscript{(59)} ہے ؛ چچا کے پاس کالا
\VUgras{لمبا کوٹ} \textsuperscript{(57)} اور سفید \VUgras{واسکٹ} ہیں۔
یہ رہے پھر \VUgras{چچا زاد بھائی} \textsuperscript{(60)} اور
\VUgras{چچا زاد بہن} \textsuperscript{(61)}۔ یہ نومولود کی
\VUgras{بہن} \textsuperscript{(62)}، ننھی برتھا کا، ہاتھ پکڑے ہوئے
ہے۔

%%K t03-c1-20-1 p
20. --- برتھا کا دوسرا \VUgras{بھائی} \textsuperscript{(63)} پیچھے چل
رہا ہے۔ وہ ایک \VUgras{رشتے دار} \textsuperscript{(64)} کو ہاتھ دیتا
ہے۔ خاندان کے \VUgras{دوست} \textsuperscript{(65)} بعد میں آتے ہیں۔

%%K t03-tit-7 sub
\VUsaut{4.6mm}
\VUpk{10.2pt}{40}{دیکھنے والے۔}
\VUsaut{3.4mm}

%%K t03-c1-21-1 p
21. --- جلوس کے قریب، یہ رہا ایک \VUgras{فقیر} \textsuperscript{(66)} ؛
وہ اپنی \VUgras{بیساکھی} \textsuperscript{(67)} پر ٹیک لگائے ہوئے ہے۔
وہ خیرات مانگ رہا ہے۔

%%K t03-c1-22-1 p
22. --- دوسری طرف ایک بہت \VUgras{مہذب} \textsuperscript{(68)} لڑکا
ہے۔ وہ اپنے جان پہچان والوں کو سلام کرتا اور « \VUgras{السلام
علیکم} » کہتا ہے۔

%%K t03-c1-23-1 p
23. --- بپتسمے کا جلوس دیکھنے کے لیے گاؤں والے اپنے گھروں سے نکل آئے
ہیں۔ ہمیں ایک \VUgras{دیہاتی} \textsuperscript{(69)} نظر آتا ہے ؛ اُس
کے سر پر \VUgras{سوتی ٹوپی} ہے، اور اُس کے پاس ایک \VUgras{دیہاتی
عورت} \textsuperscript{(70)}۔ پھر ایک \VUgras{بڑھیا}
\textsuperscript{(71)} ؛ وہ اپنی \VUgras{لاٹھی} \textsuperscript{(72)}
پر ٹیک لگائے ہوئے ہے۔ آخر میں \VUgras{چکی والا}
\textsuperscript{(73)} ؛ اُس کے پاس بڑی \VUgras{نمدے کی ٹوپی}
\textsuperscript{(74)} ہے۔ \VUgras{لکڑی کے کھڑاوے}
\textsuperscript{(75)} پہنے ایک جوان دیہاتی عورت بھی ہے ؛ وہ اپنے بچے
کو چلنا سکھا رہی ہے۔

%%K t03-c1-24-1 p
24. --- یہ منظر بہت مزے کا ہے۔

\VUsaut{4.0mm}
%%K t03-tit-8 sub
\VUcentre{12.0pt}{0}{\textit{دوسرا منظر۔}}
\VUsaut{3.0mm}

%%K t03-tit-9 sub
\VUpk{10.2pt}{40}{عام میلہ۔}
\VUsaut{2.4mm}

%%K t03-tit-10 sub
\VUpk{10.2pt}{40}{پانی کے کھیل اور کشتی دوڑ۔}
\VUsaut{3.0mm}

%%K t03-c2-01-1 p
1. --- \VUgras{گرمی} آ گئی۔ \VUgras{جون} کے مہینے کی (جولائی یا
\VUgras{اگست}) تیز دھوپ نوجوانوں کو نہانے اور کشتی چلانے پر اکساتی
ہے۔

%%K t03-c2-02-1 p
2. --- دریا کے دائیں کنارے پر کشتی ران پانی کے کھیلوں اور کشتی دوڑ
میں شریک ہونے کے لیے کپڑے اتار رہے ہیں۔

%%K t03-c2-03-1 p
3. --- اُن میں سے ایک اپنے \VUgras{جوتے} \textsuperscript{(1)} اتار
رہا ہے ؛ اُن کے تسمے کھول رہا ہے (بٹن کھول رہا ہے)۔ اُس کے دو ساتھی
پہلے ہی تیار ہیں۔ اب اُن کے پاس صرف اپنی \VUgras{بُنی ہوئی بنیان}
\textsuperscript{(2)} اور \VUgras{نہانے کا جانگیا}
\textsuperscript{(3)} ہیں۔ اپنے دوستوں کا انتظار کرتے ہوئے باتیں کر
رہے ہیں۔ چھٹی کی اور دوسرے کنارے پر ہونے والے میلے کی بات کر رہے
ہیں۔

%%K t03-c2-04-1 p
4. --- اُن کے قریب زمین پر اُن کے ساتھیوں کے \VUgras{کپڑے} پڑے ہیں :
ایک جوڑا \VUgras{چھوٹے بوٹ} \textsuperscript{(4)}، \VUgras{جوتے}
\textsuperscript{(5)}، \VUgras{جرابیں} \textsuperscript{(6)}،
\VUgras{نمدے} \textsuperscript{(7)} اور \VUgras{تنکے}
\textsuperscript{(8)} کی ٹوپیاں، اور ایک \VUgras{گلوبند}
\textsuperscript{(9)}۔

%%K t03-c2-05-1 p
5. --- نوجوانوں میں سے ایک نے اپنا \VUgras{کوٹ} \textsuperscript{(10)}
احتیاط سے تہ کیا۔ اُسے ایک تولیے پر رکھ رہا ہے ؛ اُس میں اُس کا ساتھی
جلد ہی دونوں \VUgras{کپڑے} رکھ دے گا۔

%%K t03-c2-06-1 p
6. --- چھٹا لڑکا دیر کر رہا ہے۔ اُس کی \VUgras{ٹوپی}
\textsuperscript{(11)} ابھی سر پر ہے۔ اُس نے نہ اپنی \VUgras{قمیص}
\textsuperscript{(12)} اتاری، نہ اپنا \VUgras{کالر} \textsuperscript{(13)}،
نہ اپنے \VUgras{کف} \textsuperscript{(14)}۔ ہاتھ میں اپنی
\VUgras{واسکٹ} \textsuperscript{(17)} پکڑے ہوئے ہے ؛ اُس کا
\VUgras{پتلون} \textsuperscript{(16)} اور \VUgras{گیلس}
\textsuperscript{(15)} ابھی موجود ہیں۔ اگلی دوڑ کے لیے وہ یقیناً تیار
نہ ہو گا۔

%%K t03-c2-07-1 p
7. --- نوجوانوں کے قریب ایک پگڈنڈی دریا کے کنارے تک جاتی ہے ؛ اُس کے
دونوں طرف بہت سی \VUgras{کانٹے دار جھاڑیاں} \textsuperscript{(18)}
ہیں۔ وہاں یاکوبس بابا \VUgras{سرکنڈوں} \textsuperscript{(19)} کے بیچ
\VUgras{آتش بازی} \textsuperscript{(20)} تیار کر رہے ہیں ؛ شام کو وہ
چھوڑی جائے گی۔

%%K t03-tit-11 sub
\VUsaut{4.4mm}
\VUpk{10.2pt}{40}{دوڑ۔}
\VUsaut{3.4mm}

%%K t03-c2-08-1 p
8. --- دریا پر کچھ عورتیں اپنی چھتریوں سے خود کو چھپائے
\VUgras{کشتی} \textsuperscript{(21)} میں سیر کر رہی ہیں۔ وہ بہت دلچسپ
دوڑ دیکھ رہی ہیں۔ ہر \VUgras{چھوٹی کشتی} \textsuperscript{(22)} میں
چار \VUgras{چپو والے} \textsuperscript{(24)} اپنی پوری طاقت سے چپو چلا
رہے ہیں ؛ اُسی وقت \VUgras{پتوار پکڑنے والا} \textsuperscript{(26)}
\VUgras{پتوار} \textsuperscript{(25)} تھامے اُس نازک کشتی کو چلا رہا
ہے۔ \VUgras{چپو} \textsuperscript{(23)} تال کے ساتھ پانی پر پڑتے ہیں ؛
ہلکی کشتیاں چکرا دینے والی تیزی سے چل رہی ہیں۔

%%K t03-c2-09-1 p
9. --- کچھ نوجوان پُل کے ستون کے قریب \VUgras{کشتی کے آگے والے کھمبے}
\textsuperscript{(28)} کے انعام کے لیے مقابلہ کر رہے ہیں۔ اُن میں سے
ایک، سرے پر بندھے چھوٹے \VUgras{جھنڈے} \textsuperscript{(29)} تک
پہنچنے کے لیے، اُس لچکدار اور پھسلنے والے کھمبے پر احتیاط سے چل رہا
ہے۔ اُس کا بدقسمت \VUgras{مقابل} \textsuperscript{(30)} پانی میں گر
گیا۔ وہ کنارے تک پہنچنے کے لیے تیر رہا ہے۔

%%K t03-c2-10-1 p
10. --- بڑی عمر کے نوجوان \VUgras{کشتیوں کی لڑائی}
\textsuperscript{(32)} کر رہے ہیں ؛ وہ \VUgras{پکے گھاٹ}
\textsuperscript{(33)} کے قریب ہوتی ہے۔ اِن سب کھیلوں کو ایک بہت بڑا
\VUgras{مجمع} \textsuperscript{(31)} دیکھ رہا ہے ؛ وہ \VUgras{پُل}
\textsuperscript{(27)} کی جنگلے دار دیوار پر ٹیک لگائے اور
\VUgras{کنارے} پر جمع ہے۔ پھر بھی کچھ تماشائی \VUgras{انعام کے
کھمبے} \textsuperscript{(45)} کا کھیل دیکھنے یا \VUgras{انعام بانٹنے}
کے لیے آنے والے \VUgras{شہری جلوس} کو سراہنے کے لیے مڑ کر دیکھتے
ہیں۔

%%K t03-c2-11-1 p
11. --- سب سے پہلے، یہ رہے کچھ \VUgras{لڑکے} \textsuperscript{(35)} ؛
وہ \VUgras{موسیقاروں} \textsuperscript{(36)} کے آگے چل رہے ہیں۔ پھر
\VUgras{میئر} \textsuperscript{(37)}، نائب میئر اور قصبے کی
\VUgras{بلدیہ} آتے ہیں۔ آخر میں \VUgras{آتش بجھانے والے}
\textsuperscript{(38)} چل رہے ہیں۔

%%K t03-tit-12 sub
\VUsaut{5.0mm}
\VUpk{10.2pt}{40}{میلہ۔}
\VUsaut{3.6mm}

%%K t03-c2-12-1 p
12. --- \VUgras{میلے کے میدان} \textsuperscript{(39)} میں بہت سے لوگ
ہیں۔ طرح طرح کے \VUgras{کھوکھے} دیکھنے آتے ہیں : \VUgras{لکڑی کے
گھوڑے} \textsuperscript{(40)} ؛ \VUgras{سرکس} \textsuperscript{(41)}،
جہاں تماشا شروع بھی ہو چکا ہے ؛ \VUgras{عام ناچ}
\textsuperscript{(42)}، جہاں شہر کے نوجوان لڑکے اور لڑکیاں
\VUgras{ناچنے} کا لطف اٹھاتے ہیں۔

%%K t03-c2-13-1 p
13. --- پیچھے کی طرف ہمیں \VUgras{بیل کا اکھاڑا} \textsuperscript{(44)}
نظر آتا ہے ؛ وہاں بہادر ہسپانوی \VUgras{بیل باز} غصے بھرے
\VUgras{بیل} \textsuperscript{(45)} سے لڑتا ہے۔ پھر \VUgras{پھسلن
پٹری} \textsuperscript{(49)}، \VUgras{نشانہ بازی کی جگہ}
\textsuperscript{(46)} — وہاں \VUgras{بندوقیں} چلانے اور
\VUgras{نشانے} پر مارنے کی مشق کی جاتی ہے۔

%%K t03-c2-14-1 p
14. --- قریب ہی، یہ رہا \VUgras{میلے کا تماشا گھر}
\textsuperscript{(47)} ؛ وہاں \VUgras{مزاحیہ} اور \VUgras{المیہ} کھیل
کھیلے جاتے ہیں۔ بالکل پاس \VUgras{بہت بڑے آدمی} \textsuperscript{(48)}
کا کھوکھا اور \VUgras{بہت چھوٹے آدمی} کا کھوکھا ہیں۔ پہلے کا
\VUgras{قد} دو میٹر پندرہ سنٹی میٹر ہے ؛ دوسرا ایک بچے جتنا بھی
نہیں۔

%%K t03-c2-15-1 p
15. --- آخر میں، یہ رہا بڑا \VUgras{چڑیا گھر} \textsuperscript{(50)} ؛
اُس میں \VUgras{جنگلی جانور} ہیں — زور دار \VUgras{پنجوں} والا
\VUgras{شیر} \textsuperscript{(51)}، گھنے \VUgras{ایال} والا
\VUgras{ببر شیر} \textsuperscript{(52)}، دو \VUgras{زرافے}
\textsuperscript{(53)}، اور اپنی \VUgras{سونڈ} کو بہت مہارت سے
استعمال کرنے والا \VUgras{ہاتھی} \textsuperscript{(54)}۔

%%K t03-c2-16-1 p
16. --- اُن جانوروں کو سدھانے والا \VUgras{سدھانے والا}
\textsuperscript{(55)} کھوکھے کے دروازے پر ہے۔

\VUsaut{6.0mm}
\VUfilet{22mm}
